\subsection{问题三：固定尺寸下热质耦合的“干燥终点”模型}
\subsubsection{问题三模型的建立}
问题三沿用式\eqref{eq:q2pde}，半径仍为$R_0=0.02$ m。附件1仅覆盖4 h，但模型预期运行两三天，因此必须对数据之外的边界输入进行闭合。观察后期输入已接近稳定，取3--4 h共61个观测点的算术均值：
\begin{equation}
T_a(t)=49.9989344262\ ^\circ\mathrm C,\qquad
C_e(t)=0.0499875410,\qquad t>14400\ \mathrm s .
\label{eq:plateau}
\end{equation}

该方法不把最后单个测量值的波动永久延伸到全部后期。4 h前保持线性插值；4 h处分段结束积分，以当时内部状态作为后一段初值，避免跨过输入规则的微小跳变。内部温湿状态连续，外部输入可以在分段点按指定规则切换。

按照“先看中心、再查全域”的顺序，首先定义中心含水率及候选临界时刻
\begin{equation}
C_c(t)=C(0,t),\qquad
t_c=\inf\{t\ge0:C_c(t)<0.15\}.
\end{equation}

中心是径向扩散路径最深的位置，适合作为预判点，但仅有中心值尚不能独立证明全域达标。对候选时刻扫描整个半径，定义最终条件
\begin{equation}
M(t)=\max_{0\le r\le R_0}C(r,t),\qquad
t_*=\inf\{t\ge0:M(t)<0.15\}.
\label{eq:event}
\end{equation}

若中心穿越时所有位置均不高于0.15，则确认$t_*=t_c$；若还有超标位置，中心预判不被接受，继续推进并复核。等号对应达标临界，严格低于要求发生在穿越之后，本文四位小数时长报告临界值。

全域复核不能以平均值替代。平均含水率
\begin{equation}
\overline C(t)=
\frac{\int_V\rho_d C\dd V}{\int_V\rho_d\dd V}
=\frac{2}{R_0^2}\int_0^{R_0}C(r,t)r\dd r
\end{equation}

\subsubsection{问题三模型的求解}
对离散场先考察中心节点$C_0(t)$，其下降穿越给出候选时刻；随后计算全部节点的最大值与超标集合：
\begin{equation}
g_c(t)=C_0(t)-0.15,\qquad
\mathcal E(t)=\{i:C_i(t)>0.15\}.
\end{equation}

候选时刻按连续解定位，而不对齐60 s输出点。在$t_c$处从$i=0$扫描至$N$，只有$\mathcal E(t_c)$为空时才确认全域达标临界；若非空，继续推进直到所有位置通过复核。实际严格低于时，对全部节点使用$C_i<0.15$检查，不能把等号或四位小数的显示当作已严格干燥。

模型的临界时刻为
\begin{equation}
t_*^{(3)}=206906.3987685\ \mathrm s
=57.4739996579\ \mathrm h .
\end{equation}

\begin{table}[H]\centering
\caption{问题3：固定半径下全程干基含水率（kg/kg）}\label{tab:q3}
\begin{tabular}{cccccc}\toprule
时间/h & 0 cm & 0.5 cm & 1 cm & 1.5 cm & 2 cm\\ \midrule
6.0000 & 1.0170 & 0.9881 & 0.9015 & 0.7550 & 0.5328\\
12.0000 & 0.4561 & 0.4432 & 0.4031 & 0.3297 & 0.1642\\
18.0000 & 0.2988 & 0.2916 & 0.2687 & 0.2253 & 0.0845\\
24.0000 & 0.2378 & 0.2327 & 0.2167 & 0.1855 & 0.0665\\
30.0000 & 0.2058 & 0.2018 & 0.1892 & 0.1641 & 0.0598\\
36.0000 & 0.1859 & 0.1825 & 0.1718 & 0.1504 & 0.0566\\
42.0000 & 0.1721 & 0.1691 & 0.1597 & 0.1407 & 0.0548\\
48.0000 & 0.1618 & 0.1592 & 0.1507 & 0.1334 & 0.0537\\
54.0000 & 0.1539 & 0.1514 & 0.1436 & 0.1277 & 0.0530\\
57.4740 & 0.1500 & 0.1477 & 0.1402 & 0.1249 & 0.0526\\
\bottomrule\end{tabular}\end{table}


表\ref{tab:q3}给出每6 h的五个径向位置及临界终点。12 h时中心含水率为0.4561 kg/kg，表面为0.1642 kg/kg；18 h时表面已低至0.0845 kg/kg，而中心仍为0.2988 kg/kg。由此可直接看出，若以表面含水率作为停止会遗漏仍偏湿的内部区域。

48 h时，距中心1 cm处含水率为0.1507 kg/kg，已接近阈值，但中心为0.1618 kg/kg。54 h时中心仍为0.1539 kg/kg，因此干燥还需持续约3.4740 h。临界终点处，中心为0.1500，其他四个展示位置均低于0.15，表面已降至0.0526 kg/kg。

临界表值保留四位小数，不能用显示为0.1500判断原始值严格大于、等于或小于阈值。连续事件定位使用未舍入的状态；结果表在计算完成后才舍入。该顺序保证“判断精度”不被“显示精度”替代。

既有\texttt{result3.xlsx}按60 s间隔保存各位置含水率，并增加临界时刻末行。每个时间点对应0--2 cm的21个位置，满足正文表格之外的空间明细需求。